\chapter{Connections}
\label{chap:connections}

Connections on a manifold $M$ modeled on a locally convex space $E$, given in charts by their
Christoffel symbols: the covariant derivative of vector fields and of lifts of curves, the torsion
and the curvature, connection forms and Cartan's structure equations, and geodesics as the curves
with zero covariant acceleration. All the objects are defined chart by chart, subject to a change
of charts law, and the main point in each case is that the resulting object does not depend on
the chart. We follow Lang \cite{Lan99} (\S\S IV.3, VIII.1--3, IX.1) and Chapter~6 of
\cite{Lar24}.

The last section is closer to \cite{Lar24}: a smooth field of projectors onto the tangent spaces
of an embedded submanifold of a locally convex space defines a symmetric connection, whose
covariant derivative is obtained by differentiating in the ambient space and projecting, and
whose geodesics are the curves with normal acceleration.

There is no good topology on spaces of continuous linear operators, so bilinear families such as
$x \mapsto \Gamma_x$ or $p \mapsto E_p$ are encoded as maps of several variables, linear in some of
them, with joint smoothness. Existence of geodesics is only proved on Banach manifolds, in
Chapter~\ref{chap:geodesics}.

\section{Connections and the covariant derivative}

Let $\tau$ be a $\Ck{2}$ map between open subsets of locally convex spaces. We write
$\dd\tau(y)(a)$ for its derivative and $\dd^{(2)}\tau(y; a, b)$ for the derivative of
$y \mapsto \dd\tau(y)(b)$ in the direction $a$; by the symmetry of second derivatives the latter is
symmetric in $(a, b)$. For two charts $i, j$ of the atlas of $M$ we write $\tau = j \circ i^{-1}$
for the transition map, defined on $i(U_i \cap U_j)$.

\begin{definition}[Connection]
  \label{def:connection}
  \lean{LocallyConvex.Manifolds.Connection}
  \leanok
  \uses{def:mb-manifold, def:mb-groupoid}
  A \emph{connection} on $M$ assigns to each chart $i$ of the atlas a family of maps
  $\Gamma_i(x) \colon E \times E \to E$, $x \in i(U_i)$, such that
  \begin{enumerate}
    \item each $\Gamma_i(x)$ is bilinear, and $(x, a, b) \mapsto \Gamma_i(x)(a, b)$ is smooth on
    $i(U_i) \times E \times E$;
    \item (change of charts law) for every pair of charts $i, j$, every $y$ in the domain of
    $\tau = j \circ i^{-1}$ and all $a, b \in E$,
    \[
      \Gamma_j(\tau y)\bigl(\dd\tau(y)(a), \dd\tau(y)(b)\bigr)
        = \dd\tau(y)\bigl(\Gamma_i(y)(a, b)\bigr) + \dd^{(2)}\tau(y; a, b).
    \]
  \end{enumerate}
  The maps $\Gamma_i$ are the \emph{Christoffel symbols} of the connection in the chart $i$.
\end{definition}

The symbols are not tensorial: they transform affinely, the defect being the second derivative of
the transition map. Our convention is that the first argument of $\Gamma$ is the direction in
which the base point moves. Continuity of each bilinear map $\Gamma_i(x)$ is a consequence of the
joint smoothness, and no topology on the space of bilinear maps is needed.

\begin{definition}[Local covariant derivative]
  \label{def:conn-local-cov-deriv}
  \lean{LocallyConvex.Manifolds.localCovDeriv}
  \leanok
  Let $\Gamma$ be a family of bilinear maps on an open set $U \subseteq E$ and let
  $X, Y \colon U \to E$. The \emph{local covariant derivative} is
  \[
    (\nabla_X Y)(y) = \dd Y(y)\bigl(X(y)\bigr) - \Gamma(y)\bigl(X(y), Y(y)\bigr).
  \]
\end{definition}

\begin{lemma}[Change of coordinates]
  \label{lem:conn-local-transition}
  \lean{LocallyConvex.Manifolds.dMB_comp_transition, LocallyConvex.Manifolds.localCovDeriv_transition}
  \leanok
  \uses{def:conn-local-cov-deriv}
  Let $\tau$ be $\Ck{2}$, and let $X_i, Y_i$ and $X_j, Y_j$ be related by $\tau$, that is
  $X_j \circ \tau = \dd\tau \cdot X_i$ and $Y_j \circ \tau = \dd\tau \cdot Y_i$, with $Y_i, Y_j$
  of class $\Ck{1}$. If $\Gamma_i, \Gamma_j$ obey the change of charts law of
  Definition~\ref{def:connection}, then
  \[
    (\nabla_{X_j} Y_j)(\tau y) = \dd\tau(y)\bigl((\nabla_{X_i} Y_i)(y)\bigr).
  \]
\end{lemma}

\begin{proof}
  \leanok
  Differentiating $Y_j \circ \tau = \dd\tau \cdot Y_i$ in the direction $w$ gives, by the
  second-order chain rule,
  $\dd Y_j(\tau y)(\dd\tau(y) w) = \dd^{(2)}\tau(y; w, Y_i y) + \dd\tau(y)(\dd Y_i(y) w)$. Take
  $w = X_i(y)$: the term $\dd^{(2)}\tau(y; X_i y, Y_i y)$ cancels with the one in the change of
  charts law of $\Gamma$, and the rest is linear in $\dd\tau(y)$.
\end{proof}

\begin{definition}[Covariant derivative]
  \label{def:conn-cov-deriv}
  \lean{LocallyConvex.Manifolds.Connection.covDeriv, LocallyConvex.Manifolds.Connection.isVectorField_covDeriv}
  \leanok
  \uses{def:connection, def:conn-local-cov-deriv, def:vector-field}
  Let $C$ be a connection on $M$ and $X, Y$ vector fields. The \emph{covariant derivative}
  $\nabla_X Y$ is the vector field whose value at $x$ is the local covariant derivative of the
  local representatives $\tilde X, \tilde Y$ of $X, Y$, computed in the preferred chart at $x$:
  $\dd\tilde Y \cdot \tilde X - \Gamma(\tilde X, \tilde Y)$.
\end{definition}

\begin{theorem}[Chart independence]
  \label{thm:conn-cov-deriv-chart}
  \lean{LocallyConvex.Manifolds.Connection.localRep_covDeriv}
  \leanok
  \uses{def:conn-cov-deriv, lem:conn-local-transition, def:tangent-bundle}
  Let $Y$ be a $\Ck{1}$ vector field and $X$ any vector field. In every chart $j$ of the atlas,
  the local representative of $\nabla_X Y$ is
  $\dd\tilde Y_j \cdot \tilde X_j - \Gamma_j(\tilde X_j, \tilde Y_j)$.
\end{theorem}

\begin{proof}
  \leanok
  Apply Lemma~\ref{lem:conn-local-transition} to the transition map from the preferred chart to
  $j$, which is smooth; the local representatives of vector fields are related by its
  derivative.
\end{proof}

So the choice of the preferred chart in Definition~\ref{def:conn-cov-deriv} is only a matter of
convenience. The usual algebraic rules follow at once from the local formula.

\begin{proposition}[Rules of the covariant derivative]
  \label{prop:conn-rules}
  \lean{LocallyConvex.Manifolds.Connection.covDeriv_smulField_left, LocallyConvex.Manifolds.Connection.covDeriv_smulField_right, LocallyConvex.Manifolds.Connection.contMDiffMB_covDeriv}
  \leanok
  \uses{thm:conn-cov-deriv-chart, def:contmdiff}
  Let $X, Y$ be vector fields with $Y$ of class $\Ck{1}$, and let $\varphi \colon M \to \R$.
  \begin{enumerate}
    \item $\nabla_{\varphi X} Y = \varphi\, \nabla_X Y$.
    \item If $\varphi$ is $\Ck{1}$, then $\nabla_X(\varphi Y) = X(\varphi)\, Y + \varphi\,\nabla_X Y$
    (Leibniz rule).
    \item If $X$ is $\Ck{k}$ and $Y$ is $\Ck{k+1}$, then $\nabla_X Y$ is $\Ck{k}$.
  \end{enumerate}
\end{proposition}

\begin{proof}
  \leanok
  The first two items are computed in the preferred chart: $\Gamma$ is linear in its first
  argument, and the derivative of $\varphi \tilde Y$ in the direction $\tilde X$ is
  $\dd\varphi(\tilde X)\,\tilde Y + \varphi\,\dd\tilde Y(\tilde X)$. For the third, in each chart
  $y \mapsto \dd\tilde Y(y)(\tilde X y)$ and $y \mapsto \Gamma(y)(\tilde X y, \tilde Y y)$ are
  compositions of $\Ck{k}$ maps.
\end{proof}

\begin{proposition}[Torsion]
  \label{prop:conn-torsion}
  \lean{LocallyConvex.Manifolds.Connection.covDeriv_sub_covDeriv_sub_lieBracket, LocallyConvex.Manifolds.Connection.IsSymmetric, LocallyConvex.Manifolds.Connection.covDeriv_sub_covDeriv}
  \leanok
  \uses{def:conn-cov-deriv}
  For vector fields $X, Y$, the vector $\nabla_X Y - \nabla_Y X - [X, Y]$ at $x$ is, in the
  preferred chart, $\Gamma(\tilde Y, \tilde X) - \Gamma(\tilde X, \tilde Y)$; in particular it only
  depends on $X(x)$ and $Y(x)$. We say that the connection is \emph{symmetric} (torsion free) if
  its symbols are symmetric in every chart; in that case $\nabla_X Y - \nabla_Y X = [X, Y]$.
\end{proposition}

\begin{proof}
  \leanok
  The Lie bracket in a chart is $\dd\tilde Y \cdot \tilde X - \dd\tilde X \cdot \tilde Y$, and
  the derivative terms cancel.
\end{proof}

\begin{remark}
  Symmetry of a connection is a chart-free condition: the change of charts law shows that if
  $\Gamma_i$ is symmetric then so is $\Gamma_j$, because $\dd^{(2)}\tau$ is symmetric and
  $\dd\tau(y)$ is invertible on the overlap. In the Lean code the definition quantifies over all
  the charts of the atlas, which avoids this argument.
\end{remark}

\section{Curvature}

We use Lang's sign for the curvature.

\begin{definition}[Curvature]
  \label{def:conn-curvature}
  \lean{LocallyConvex.Manifolds.dSymb, LocallyConvex.Manifolds.localCurv, LocallyConvex.Manifolds.Connection.curvature}
  \leanok
  \uses{def:conn-cov-deriv}
  For vector fields $X, Y, Z$, the \emph{curvature} at $x$ is
  \[
    R(X, Y) Z = \nabla_X \nabla_Y Z - \nabla_Y \nabla_X Z - \nabla_{[X, Y]} Z .
  \]
  Locally, for a family of symbols $\Gamma$ on an open set of $E$, we write
  $\dd\Gamma_x(v)(a, b)$ for the derivative of $x \mapsto \Gamma(x)(a, b)$ in the direction $v$,
  and
  \[
    R_x(u, v) w = \dd\Gamma_x(v)(u, w) - \dd\Gamma_x(u)(v, w)
      + \Gamma_x\bigl(u, \Gamma_x(v, w)\bigr) - \Gamma_x\bigl(v, \Gamma_x(u, w)\bigr).
  \]
\end{definition}

\begin{lemma}[Local curvature formula]
  \label{lem:conn-local-curv}
  \lean{LocallyConvex.Manifolds.dMB_symb_comp, LocallyConvex.Manifolds.localCovDeriv_curv}
  \leanok
  \uses{def:conn-curvature, def:conn-local-cov-deriv}
  Let $X, Y, Z$ be $\Ck{2}$ fields on an open set $U \subseteq E$ and $\Gamma$ a $\Ck{1}$ family of
  bilinear maps. Then, pointwise on $U$,
  $\nabla_X \nabla_Y Z - \nabla_Y \nabla_X Z - \nabla_{[X, Y]} Z = R(X, Y) Z$
  with $R$ the trilinear expression of Definition~\ref{def:conn-curvature}.
\end{lemma}

\begin{proof}
  \leanok
  Expand both covariant derivatives. The derivative of $y \mapsto \Gamma(y)(A y, B y)$ is
  $\dd\Gamma_y(h)(A y, B y) + \Gamma_y(\dd A(y) h, B y) + \Gamma_y(A y, \dd B(y) h)$. All the
  terms containing second derivatives of $Z$ cancel by the symmetry of second derivatives, and
  the terms containing first derivatives of $X$, $Y$, $Z$ cancel with $\nabla_{[X,Y]} Z$ and
  among themselves. This is Lang \cite{Lan99}, Prop.~IX.1.2.
\end{proof}

\begin{theorem}[The curvature is tensorial]
  \label{thm:conn-curvature-tensorial}
  \lean{LocallyConvex.Manifolds.Connection.curvature_eq}
  \leanok
  \uses{lem:conn-local-curv, thm:conn-cov-deriv-chart}
  Let $X, Y, Z$ be $\Ck{2}$ vector fields. Then $R(X, Y) Z$ at $x$ is the trilinear expression
  $R_{\tilde x}(\tilde X, \tilde Y)\tilde Z$ in the preferred chart, evaluated at the values of
  $X, Y, Z$ at $x$. In particular, it only depends on $X(x)$, $Y(x)$ and $Z(x)$.
\end{theorem}

\begin{proof}
  \leanok
  Near $x$ the local representatives of $\nabla_Y Z$, $\nabla_X Z$ and $[X, Y]$ in the preferred
  chart at $x$ are given by the local formulas (Theorem~\ref{thm:conn-cov-deriv-chart}), and the
  local covariant derivative at a point only depends on the germ of the second field. Then apply
  Lemma~\ref{lem:conn-local-curv}. Compare \cite{Lan99}, Prop.~IX.1.1--1.2.
\end{proof}

\begin{proposition}[Symmetries of the curvature]
  \label{prop:conn-curvature-symmetries}
  \lean{LocallyConvex.Manifolds.localCurv_swap, LocallyConvex.Manifolds.localCurv_bianchi, LocallyConvex.Manifolds.Connection.curvature_swap, LocallyConvex.Manifolds.Connection.curvature_bianchi}
  \leanok
  \uses{thm:conn-curvature-tensorial, prop:conn-torsion}
  For $\Ck{2}$ vector fields, $R(Y, X) Z = -R(X, Y) Z$. If the connection is symmetric, the first
  Bianchi identity holds: $R(X, Y) Z + R(Y, Z) X + R(Z, X) Y = 0$.
\end{proposition}

\begin{proof}
  \leanok
  Both identities are checked on the trilinear expression. Skew-symmetry is immediate. For
  Bianchi, if $\Gamma$ is symmetric then so is $\dd\Gamma_x(v)$, and the twelve terms of the
  cyclic sum cancel in pairs.
\end{proof}

\section{Connection forms and the structure equations}
\label{sec:connection-forms}

Classically, the connection form of a connection on a vector bundle is the $\mathrm{End}(F)$-valued
$1$-form $A(y)(v) = \Gamma(y)(v, \cdot)$ in a local trivialization, so that $\nabla = \dd - A$, and
the curvature is $\Omega = -\dd A + A \wedge A$. If the fiber $F$ is not normable, no vector
topology on $L(F)$ makes the evaluation $L(F) \times F \to F$ continuous \cite[Satz 1]{Mai63}, so
$A$ is not a
differential form with values in a locally convex space of operators. We encode
$\mathrm{End}(F)$-valued forms as families of $F$-valued forms depending linearly on an extra
argument $\xi \in F$. The local statements are proved for an arbitrary fiber $F$; the tangent
bundle is the case $F = E$.

\begin{definition}[Connection and curvature forms]
  \label{def:connection-forms}
  \lean{LocallyConvex.Manifolds.IsLocalSymbol, LocallyConvex.Manifolds.symbD,
    LocallyConvex.Manifolds.bundleCurv, LocallyConvex.Manifolds.connForm,
    LocallyConvex.Manifolds.connWedge, LocallyConvex.Manifolds.curvLoc,
    LocallyConvex.Manifolds.IsLocalSymbol.isLocalForm_connForm}
  \leanok
  \uses{def:local-form}
  Let $U \subseteq E$ be open and $\Gamma \colon U \times E \times F \to F$ smooth and bilinear in
  the last two variables. For $\xi \in F$, the \emph{connection form} is the local $1$-form
  $A^\xi(y)(v) = \Gamma(y)(v, \xi)$, and the \emph{curvature form} is the local $2$-form
  \[
    \Omega^\xi(y)(u, v) = \partial_v\Gamma(y)(u, \xi) - \partial_u\Gamma(y)(v, \xi)
      + \Gamma(y)(u, \Gamma(y)(v, \xi)) - \Gamma(y)(v, \Gamma(y)(u, \xi)).
  \]
  We write $(A \wedge A)(u, v)\xi = A(u)A(v)\xi - A(v)A(u)\xi$.
\end{definition}

\begin{theorem}[Cartan's structure equation]
  \label{thm:structure-equation}
  \lean{LocallyConvex.Manifolds.dLoc_connForm, LocallyConvex.Manifolds.curvLoc_eq,
    LocallyConvex.Manifolds.bundleCurv_eq_localCurv,
    LocallyConvex.Manifolds.Connection.curvature_eq_structure}
  \leanok
  \uses{def:connection-forms, def:local-d, thm:conn-curvature-tensorial}
  $\Omega^\xi = -\dd A^\xi + (A \wedge A)\xi$. For the tangent bundle, $\Omega$ is the local
  curvature of Lemma~\ref{lem:conn-local-curv}, so the curvature of a connection on $M$, evaluated
  on vector fields, is computed in the preferred chart by $-\dd A + A \wedge A$.
\end{theorem}

\begin{proof}
  \leanok
  $\dd A^\xi(y)(u, v) = \partial_u\Gamma(y)(v, \xi) - \partial_v\Gamma(y)(u, \xi)$.
\end{proof}

\begin{theorem}[Second Bianchi identity]
  \label{thm:second-bianchi}
  \lean{LocallyConvex.Manifolds.IsLocalSymbol.hasDirDerivAt_comp,
    LocallyConvex.Manifolds.IsLocalSymbol.dLoc_curvLoc,
    LocallyConvex.Manifolds.Connection.dLoc_curvLoc}
  \leanok
  \uses{thm:structure-equation, thm:local-dd}
  The covariant exterior derivative of the curvature vanishes, $D\Omega = \dd\Omega - [A \wedge
  \Omega] = 0$: for $u, v, t \in E$ and $\xi \in F$,
  \[
    \dd\Omega(u, v, t)\xi = [A(u), \Omega(v, t)]\xi - [A(v), \Omega(u, t)]\xi
      + [A(t), \Omega(u, v)]\xi,
  \]
  where $[A(a), \Omega(b, c)]\xi = A(a)\Omega(b, c)\xi - \Omega(b, c)A(a)\xi$.
\end{theorem}

\begin{proof}
  \leanok
  By the structure equation and $\dd \circ \dd = 0$, $\dd\Omega^\xi = \dd((A \wedge A)\xi)$, which
  is computed with the product rule
  $\partial_c(\Gamma(a, \Gamma(b, \xi))) = \partial_c\Gamma(a, \Gamma(b, \xi)) + \Gamma(a,
  \partial_c\Gamma(b, \xi))$. The terms with three factors $\Gamma$ on the right cancel by the
  Jacobi identity for commutators.
\end{proof}

\begin{proposition}[Gauge transformations; torsion]
  \label{prop:gauge}
  \lean{LocallyConvex.Manifolds.Connection.isLocalSymbol,
    LocallyConvex.Manifolds.Connection.connForm_transition,
    LocallyConvex.Manifolds.torsionLoc,
    LocallyConvex.Manifolds.Connection.torsionLoc_transition,
    LocallyConvex.Manifolds.Connection.isSymmetric_iff_torsionLoc}
  \leanok
  \uses{def:connection, def:connection-forms, prop:conn-torsion}
  Let $\tau = j \circ i^{-1}$ and $g(y) = \dd\tau(y)$. The connection forms of a connection on
  $M$ satisfy $A_j(\tau y)(g v)(g w) = g\,A_i(y)(v) w + (\partial_v g)(y) w$, so $A$ is not a
  tensor. The torsion $T(u, v) = A(u)v - A(v)u = \Gamma(u, v) - \Gamma(v, u)$ is a tensor:
  $T_j(\tau y)(g a, g b) = g\, T_i(y)(a, b)$.
\end{proposition}

\begin{proof}
  \leanok
  The first statement is the change of charts law of the symbols. For the torsion, the
  inhomogeneous terms $\dd^2\tau(a, b)$ and $\dd^2\tau(b, a)$ cancel.
\end{proof}

\begin{remark}
  For a vector bundle with fiber $F$ the same formalism applies to the symbols in local
  trivializations, with transition functions $g(y) \in GL(F)$ in place of $\dd\tau(y)$; for
  non-normable $F$ one can only require $(y, \xi) \mapsto g(y)\xi$ to be smooth, not
  $y \mapsto g(y)$. Vector bundles, their connections and the tensorial transformation law of
  $\Omega$ are treated in Section~\ref{sec:vector-bundles}.
\end{remark}

\section{Covariant derivative along curves and geodesics}

We use the $\Ck{k}$ curves of Definition~\ref{def:ck-curve} and the velocity of
Definition~\ref{def:velocity}. Let $\alpha \colon \R \to M$ be a curve. A \emph{lift} of $\alpha$
is a map $\beta \colon \R \to TM$ with $\beta(t) \in T_{\alpha(t)} M$. In a chart $i$ we write
$\tilde\alpha = i \circ \alpha$ and $\mu_i(t)$ for the fiber coordinate of $\beta(t)$.

\begin{definition}[Lifts of curves]
  \label{def:conn-lift}
  \lean{LocallyConvex.Manifolds.liftCoord, LocallyConvex.Manifolds.IsCkLift}
  \leanok
  \uses{def:ck-curve, def:tangent-bundle}
  A lift $\beta$ of $\alpha$ is $\Ck{k}$ if, for every chart $i$, its fiber coordinate $\mu_i$ is
  $\Ck{k}$ on $\alpha^{-1}(U_i)$.
\end{definition}

\begin{lemma}[Change of charts along a curve]
  \label{lem:conn-curve-transition}
  \lean{LocallyConvex.Manifolds.tangentChart_snd_transition, LocallyConvex.Manifolds.dMB_dMB_curve}
  \leanok
  \uses{def:conn-lift, def:velocity}
  For two charts $i, j$ and $\tau = j \circ i^{-1}$, we have
  $\mu_j(t) = \dd\tau(\tilde\alpha_i t)(\mu_i t)$ and
  $\dot{\tilde\alpha}_j = \dd\tau(\tilde\alpha_i)(\dot{\tilde\alpha}_i)$, and
  \[
    \dot\mu_j = \dd^{(2)}\tau\bigl(\tilde\alpha_i; \dot{\tilde\alpha}_i, \mu_i\bigr)
      + \dd\tau(\tilde\alpha_i)(\dot\mu_i).
  \]
  The fiber coordinate of the velocity in any chart $j$ is $\dot{\tilde\alpha}_j$, and the
  velocity of a $\Ck{k+1}$ curve is a $\Ck{k}$ lift.
\end{lemma}

\begin{proof}
  \leanok
  The first identity is the change of charts of the tangent bundle, and the others follow from
  the chain rule and the product rule applied to $t \mapsto \dd\tau(\tilde\alpha_i t)(\mu_i t)$.
\end{proof}

\begin{definition}[Covariant derivative along a curve]
  \label{def:cov-deriv-along}
  \lean{LocallyConvex.Manifolds.Connection.covDerivAlong}
  \leanok
  \uses{def:connection, def:conn-lift}
  Let $C$ be a connection and $\beta$ a lift of $\alpha$. The \emph{covariant derivative of
  $\beta$ along $\alpha$} at $t$ is, in the preferred chart at $\alpha(t)$,
  \[
    D_t \beta = \dot\mu - \Gamma(\tilde\alpha)\bigl(\dot{\tilde\alpha}, \mu\bigr).
  \]
\end{definition}

\begin{theorem}[Chart independence along curves]
  \label{thm:conn-cov-deriv-along-chart}
  \lean{LocallyConvex.Manifolds.Connection.covDerivAlong_chart}
  \leanok
  \uses{def:cov-deriv-along, lem:conn-curve-transition}
  Let $\alpha$ be a $\Ck{1}$ curve and $\beta$ a $\Ck{1}$ lift. Read in any chart $j$ around
  $\alpha(t)$, the covariant derivative $D_t\beta$ is
  $\dot\mu_j - \Gamma_j(\tilde\alpha_j)(\dot{\tilde\alpha}_j, \mu_j)$.
\end{theorem}

\begin{proof}
  \leanok
  By Lemma~\ref{lem:conn-curve-transition}, $\dot\mu_j$ contains the term
  $\dd^{(2)}\tau(\tilde\alpha_i; \dot{\tilde\alpha}_i, \mu_i)$, which cancels with the one in the
  change of charts law of $\Gamma$ applied to $(\dot{\tilde\alpha}_i, \mu_i)$. The remaining
  terms are $\dd\tau(\tilde\alpha_i)$ applied to the expression in the chart $i$. Compare
  \cite{Lan99}, \S VIII.3.
\end{proof}

\begin{definition}[Geodesic]
  \label{def:geodesic}
  \lean{LocallyConvex.Manifolds.Connection.IsGeodesic}
  \leanok
  \uses{def:cov-deriv-along, def:velocity}
  A curve $\alpha$ is a \emph{geodesic} of the connection $C$ if its velocity is parallel:
  $D_t \alpha' = 0$ for every $t$.
\end{definition}

\begin{theorem}[The geodesic equation]
  \label{thm:conn-geodesic-equation}
  \lean{LocallyConvex.Manifolds.Connection.isGeodesic_iff}
  \leanok
  \uses{def:geodesic, thm:conn-cov-deriv-along-chart}
  A $\Ck{2}$ curve $\alpha$ is a geodesic if and only if, in every chart $j$,
  \[
    \ddot{\tilde\alpha}_j = \Gamma_j(\tilde\alpha_j)\bigl(\dot{\tilde\alpha}_j, \dot{\tilde\alpha}_j\bigr).
  \]
\end{theorem}

\begin{proof}
  \leanok
  By Theorem~\ref{thm:conn-cov-deriv-along-chart} and Lemma~\ref{lem:conn-curve-transition}, in
  the chart $j$ the covariant derivative of the velocity is
  $\ddot{\tilde\alpha}_j - \Gamma_j(\dot{\tilde\alpha}_j, \dot{\tilde\alpha}_j)$, transported by
  the invertible map $\dd\tau$. For the converse it suffices to use the preferred chart.
\end{proof}

In terms of the spray of the connection (Proposition~\ref{prop:conn-to-spray}), the geodesic
equation reads $\ddot{\tilde\alpha} = f(\tilde\alpha, \dot{\tilde\alpha})$ in every chart, which is
the classical definition of the geodesics of a spray.

\begin{remark}[Existence of geodesics]
  \label{rem:conn-no-existence}
  We do not claim that geodesics exist, nor that they are unique, for given initial position and
  velocity. In a Banach space the geodesic equation is an ordinary differential equation with a
  smooth right hand side, and the Picard--Lindel\"of theorem gives local existence and
  uniqueness; this is how the exponential map is constructed in \cite{Lan99} and \cite{Lar24}.
  In a general locally convex space, and even in Fr\'echet spaces, ordinary differential equations
  with smooth right hand side may have no solutions, or many \cite[Part~I, 5.6.1]{Ham82}. Therefore existence of geodesics,
  the exponential map and parallel transport along curves are not part of the theory developed
  here; they would have to be proved case by case, or under additional hypotheses on $E$.
  What we can do in general is to recognize geodesics, as in Theorem~\ref{thm:projector-geodesics}
  below. When the model space is a Banach space, existence and uniqueness are proved in the next
  section.
\end{remark}

\section{Fields of projectors: derive and project}
\label{sec:conn-projectors}

Let $F$ be a locally convex space and $X \subseteq F$ an embedded submanifold modeled on a closed
subspace $S \subseteq F$ (Definition~\ref{def:embedded-submanifold}). Recall that in a slice chart
$\varphi$ the tangent space at $p$ is $T_p X = \{w \in F : \dd\varphi(p)(w) \in S\}$, and that it
does not depend on the slice chart (Definition~\ref{def:tangent-space-submanifold}). We write
$\psi = \varphi^{-1}|_S$ for the local parametrization of $X$ given by $\varphi$; then
$\dd\psi(y)$ maps $S$ injectively onto $T_{\psi(y)} X$.

\begin{definition}[Field of projectors]
  \label{def:projector-field}
  \lean{LocallyConvex.Manifolds.ProjectorField, LocallyConvex.Manifolds.ProjectorField.param, LocallyConvex.Manifolds.ProjectorField.localE}
  \leanok
  \uses{def:embedded-submanifold, def:tangent-space-submanifold}
  A \emph{field of projectors} onto the tangent spaces of $X$ is a map
  $E \colon X \times F \to F$, $(p, w) \mapsto E_p w$, such that
  \begin{enumerate}
    \item each $E_p$ is linear, $E_p^2 = E_p$, and the range of $E_p$ is $T_p X$;
    \item for every slice chart $\varphi$, the map $(y, w) \mapsto E_{\psi(y)} w$ is smooth on
    $\varphi(U) \cap S \times F$.
  \end{enumerate}
\end{definition}

\begin{remark}
  In \cite{Lar24} the field of projectors is a smooth map $X \to \mathcal B(F)$ into the bounded
  operators of a Banach space. For a locally convex space $F$ there is no good topology on the
  space of continuous operators, so we replace it by a map of two variables, linear in the second
  one, with joint smoothness. Note also that no complement of $S$ in $F$ is assumed: the kernels
  $\ker E_p$ are complements of the tangent spaces $T_p X$, but they are not required to come from
  a global splitting of $F$, nor to be related to each other.
\end{remark}

\begin{lemma}[Tangent vectors]
  \label{lem:projector-tangent}
  \lean{LocallyConvex.Manifolds.ProjectorField.mem_tangentSpace, LocallyConvex.Manifolds.ProjectorField.E_eq_self, LocallyConvex.Manifolds.IsEmbeddedSubmanifold.dMB_symm_mem_tangentSpace, LocallyConvex.Manifolds.ProjectorField.E_dMB_param, LocallyConvex.Manifolds.ProjectorField.dMB_param_injective}
  \leanok
  \uses{def:projector-field}
  $E_p w \in T_p X$ for every $w$, and $E_p w = w$ if $w \in T_p X$. In a slice chart,
  $\dd\psi(y)(s) \in T_{\psi(y)} X$ for $s \in S$, hence $E_{\psi(y)}\dd\psi(y)(s) = \dd\psi(y)(s)$,
  and $\dd\psi(y)$ is injective on $S$.
\end{lemma}

\begin{proof}
  \leanok
  The first two claims follow from $E_p^2 = E_p$ and the description of the range. Since
  $\dd\varphi$ is a left and right inverse of $\dd\varphi^{-1}$, we have
  $\dd\varphi(\dd\psi(y) s) = s \in S$, which gives the rest.
\end{proof}

\begin{definition}[Symbols of a field of projectors]
  \label{def:projector-symbols}
  \lean{LocallyConvex.Manifolds.ProjectorField.symb, LocallyConvex.Manifolds.ProjectorField.symb_mem, LocallyConvex.Manifolds.ProjectorField.coe_symb}
  \leanok
  \uses{def:projector-field}
  In a slice chart $\varphi$, with $\psi = \varphi^{-1}|_S$ and $y \in \varphi(U) \cap S$, put
  \[
    \Gamma_y(a, b) = -\,\dd\varphi(\psi y)\Bigl(E_{\psi(y)}\bigl(\dd^{(2)}\psi(y; a, b)\bigr)\Bigr),
    \qquad a, b \in S .
  \]
  This vector lies in $S$, because $E_{\psi(y)}(\cdot)$ is tangent.
\end{definition}

Equivalently, $\dd\psi(y)\bigl(\Gamma_y(a, b)\bigr) = -E\bigl(\dd^{(2)}\psi(y; a, b)\bigr)$: the
symbols are the tangent part of the second derivative of the parametrization, read in the chart.

\begin{lemma}[Properties of the symbols]
  \label{lem:projector-symbols}
  \lean{LocallyConvex.Manifolds.ProjectorField.symb_symm, LocallyConvex.Manifolds.ProjectorField.smoothMBOn_symb, LocallyConvex.Manifolds.ProjectorField.dMB_param_symb, LocallyConvex.Manifolds.ProjectorField.symb_transition}
  \leanok
  \uses{def:projector-symbols, lem:projector-tangent, def:connection}
  The symbols $\Gamma_y$ are bilinear and symmetric, they are smooth on the chart domain, and for
  two slice charts they obey the change of charts law of Definition~\ref{def:connection}.
\end{lemma}

\begin{proof}
  \leanok
  Symmetry and bilinearity come from those of $\dd^{(2)}\psi(y; \cdot, \cdot)$ (Schwarz), and
  smoothness from the smoothness of $E$ in slice charts and of $\psi$. For the change of charts
  law let $\tau$ be the transition map between two slice charts, so that $\psi' \circ \tau = \psi$
  near $y$. The second-order chain rule gives
  $\dd^{(2)}\psi = \dd^{(2)}\psi'(\dd\tau\, a, \dd\tau\, b) + \dd\psi'(\dd^{(2)}\tau(a, b))$. The
  last term is tangent, so $E$ fixes it; applying $-\dd\varphi' \circ E$ yields the law.
\end{proof}

\begin{theorem}[The connection of a field of projectors]
  \label{thm:projector-connection}
  \lean{LocallyConvex.Manifolds.ProjectorField.connection, LocallyConvex.Manifolds.ProjectorField.connection_isSymmetric}
  \leanok
  \uses{lem:projector-symbols, def:connection, prop:conn-torsion}
  A field of projectors $E$ onto the tangent spaces of $X$ defines a symmetric connection on $X$,
  whose symbols in the chart induced by a slice chart $\varphi$ are
  $\Gamma_y(a, b) = -\dd\varphi\bigl(E(\dd^{(2)}\psi(y; a, b))\bigr)$.
\end{theorem}

\begin{proof}
  \leanok
  The charts of the atlas of $X$ are the restrictions of slice charts, and
  Lemma~\ref{lem:projector-symbols} provides all the axioms. This is the first part of
  \cite[Teorema~6.6.1]{Lar24}.
\end{proof}

The next result explains the title of the section: the covariant derivative is obtained by
differentiating in $F$ and projecting onto the tangent space.

\begin{theorem}[Derive and project]
  \label{thm:projector-derive-project}
  \lean{LocallyConvex.Manifolds.ProjectorField.dMB_param_localCovDeriv}
  \leanok
  \uses{thm:projector-connection, def:conn-local-cov-deriv}
  Let $A, B$ be fields on an open subset of the chart domain in $S$, with $B$ of class $\Ck{1}$,
  and let $\hat B(y) = \dd\psi(y)(B y)$ be the corresponding tangent field read in $F$. Then
  \[
    \dd\psi(y)\bigl((\nabla_A B)(y)\bigr) = E_{\psi(y)}\bigl(\dd\hat B(y)(A y)\bigr).
  \]
  In other words, $\nabla_X Y = E(\dd Y \cdot X)$ when the fields are read in the ambient
  space $F$.
\end{theorem}

\begin{proof}
  \leanok
  By the second-order chain rule,
  $\dd\hat B(y)(A y) = \dd^{(2)}\psi(y; A y, B y) + \dd\psi(y)(\dd B(y)(A y))$. Apply $E$: the
  second term is tangent and is fixed, and $E(\dd^{(2)}\psi(y; A y, B y))
  = -\dd\psi(y)(\Gamma_y(A y, B y))$ by Definition~\ref{def:projector-symbols}. Compare
  \cite[\S6.6]{Lar24}.
\end{proof}

\begin{lemma}[Acceleration in a slice chart]
  \label{lem:projector-accel}
  \lean{LocallyConvex.Manifolds.ProjectorField.E_accel_eq, LocallyConvex.Manifolds.IsEmbeddedSubmanifold.isCkCurve_of_smooth}
  \leanok
  \uses{def:projector-symbols, def:ck-curve}
  Let $c$ be a $\Ck{2}$ curve in the chart domain in $S$ and $\alpha = \psi \circ c$. Then
  \[
    E_{\alpha(t)}\bigl(\ddot\alpha(t)\bigr)
      = \dd\psi(c t)\bigl(\ddot c(t) - \Gamma_{c(t)}(\dot c(t), \dot c(t))\bigr),
  \]
  where $\ddot\alpha$ is the acceleration of $\alpha$ in $F$. Moreover, a curve in $X$ which is
  smooth as a curve in $F$ is a $\Ck{k}$ curve of $X$ for every $k$.
\end{lemma}

\begin{proof}
  \leanok
  We have $\dot\alpha = \dd\psi(c)(\dot c)$ and, by the chain rule and the product rule,
  $\ddot\alpha = \dd^{(2)}\psi(c; \dot c, \dot c) + \dd\psi(c)(\ddot c)$; apply $E$ as in
  the proof of Theorem~\ref{thm:projector-derive-project}. For the last claim, the local
  representative in a slice chart is the corestriction to $S$ of $\varphi \circ \alpha$.
\end{proof}

\begin{theorem}[Derive and project: geodesics]
  \label{thm:projector-geodesics}
  \lean{LocallyConvex.Manifolds.ProjectorField.connection_isGeodesic_iff}
  \leanok
  \uses{thm:projector-connection, thm:conn-geodesic-equation, lem:projector-accel, def:geodesic}
  Let $\alpha$ be a $\Ck{2}$ curve in $X$. Then $\alpha$ is a geodesic of the connection defined
  by the field of projectors $E$ if and only if its acceleration in $F$ is normal:
  \[
    E_{\alpha(t)}\bigl(\ddot\alpha(t)\bigr) = 0 \qquad \text{for every } t.
  \]
\end{theorem}

\begin{proof}
  \leanok
  By Theorem~\ref{thm:conn-geodesic-equation}, $\alpha$ is a geodesic if and only if
  $\ddot c = \Gamma_c(\dot c, \dot c)$ in every chart, where $c$ is the local representative. By
  Lemma~\ref{lem:projector-accel} this is equivalent to $E(\ddot\alpha) = 0$, since $\dd\psi(c)$ is
  injective on $S$ (Lemma~\ref{lem:projector-tangent}). This is part of \cite[Teorema~6.6.1]{Lar24}.
\end{proof}

For instance, for the unit sphere of a (complete) Hilbert space with the orthogonal projectors
onto $p^\perp$, the geodesics are the curves in the sphere with $\ddot\alpha(t)$ proportional to
$\alpha(t)$. This example is not formalized.

Finally, we compare with the formulation of the book, where the connection is written in terms of
the derivative of the field of projectors. We denote with $E_{*p, v} w$ the derivative of
$q \mapsto E_q w$ at $p$ in the direction $v$, computed in the slice chart.

\begin{proposition}[The identity of the book]
  \label{prop:projector-book-identity}
  \lean{LocallyConvex.Manifolds.ProjectorField.dMB_localE_param, LocallyConvex.Manifolds.ProjectorField.dMB_localE_param_symm}
  \leanok
  \uses{def:projector-field, lem:projector-tangent}
  Let $v = \dd\psi(y)(a)$ and $w = \dd\psi(y)(b)$ be tangent vectors at $p = \psi(y)$, with
  $a, b \in S$. Then
  \[
    E_{*p, v}\, w = (1 - E_p)\, \dd^{(2)}\psi(y; a, b).
  \]
  In particular $E_{*p, v}\, w = E_{*p, w}\, v$.
\end{proposition}

\begin{proof}
  \leanok
  Since $\dd\psi(y')(b)$ is tangent, $E_{\psi(y')}\dd\psi(y')(b) = \dd\psi(y')(b)$ for $y'$ near
  $y$. Differentiate this identity in the direction $a$: by linearity of $E$ in the second
  variable, the left side gives $E_{*p, v} w + E_p\, \dd^{(2)}\psi(y; a, b)$, and the right side
  gives $\dd^{(2)}\psi(y; a, b)$. The symmetry follows from the symmetry of $\dd^{(2)}\psi$.
\end{proof}

\begin{remark}
  In \cite{Lar24} the connection is presented through $\Gamma_p(v, w) = E_{*p, v} w$, which by
  Proposition~\ref{prop:projector-book-identity} is the normal part $(1 - E)\dd^{(2)}\psi(a, b)$
  of the second derivative, read in $F$. The symbols of
  Theorem~\ref{thm:projector-connection} are instead the tangent part, pulled back to $S$ by the
  chart; the two descriptions are related through the decomposition
  $\dd^{(2)}\psi = E\,\dd^{(2)}\psi + (1 - E)\dd^{(2)}\psi$. In both cases the symmetry of the
  connection is a consequence of the symmetry of second derivatives of $\psi$. In the Lean code
  the connection is defined by the tangent part, and the book's formula enters only through
  Proposition~\ref{prop:projector-book-identity}.
\end{remark}
